\chapter{Relative products and all-arity gluing}
\label{v4-arith:gluing:chapter}

A relative cocycle consists of a cocycle and a primitive of its
restriction. Multiplication therefore has a choice: use the primitive
of the first input, or that of the second. The resulting products agree
on cohomology, but the chain comparison contains ordered products of
arbitrarily many primitives. These higher terms also change module
actions. They must be retained when relative classes are composed.

We construct the comparison over an arbitrary commutative ring. Its
formulas use finite sums on each tensor word and require no division.
This gives an integral comparison of relative cochains, their modules,
and their bimodules. Passing to an inverse limit of finite coefficient
rings preserves the comparison. Passing instead to an unrestricted
product over tensor lengths can destroy its continuity.

\section{The two endpoint products}
\label{v4-arith:gluing:sec:products}

Let \(R\) be a commutative ring. All complexes are cohomological, their
differentials have degree one, and tensor products use the Koszul rule.
Fix a strict cospan of associative differential graded \(R\)-algebras
\[
 A\xrightarrow{f}B\xleftarrow{g}C.
\]
The algebras may be nonunital. Define the relative complex
\begin{equation}\label{v4-arith:gluing:eq:complex}
 P^n=A^n\oplus C^n\oplus B^{n-1},\qquad
 D(a,c,y)=(d_Aa,d_Cc,f(a)-g(c)-d_By).
\end{equation}
A cocycle \((a,c,y)\) thus includes a specified primitive \(y\) of
\(f(a)-g(c)\). The identity \(D^2=0\) follows from the chain-map
identities for \(f\) and \(g\). In the derived category of complexes,
\(P\) is the fibre of \(f-g:A\oplus C\to B\). This statement specifies
its underlying complex; its multiplicative structure is the additional
construction below.

For homogeneous \(x=(a,c,y)\in P^p\) and \(x'=(a',c',y')\), put
\begin{equation}\label{v4-arith:gluing:eq:products}
 \begin{split}
 \mu_L(x,x')&=(aa',cc',(-1)^pf(a)y'+yg(c')),\\
 \mu_R(x,x')&=(aa',cc',yf(a')+(-1)^pg(c)y').
 \end{split}
\end{equation}
The symbols \(L\) and \(R\) name these displayed choices throughout the
chapter. In the relative case \(C=0\), their third components reduce to
\((-1)^pf(a)y'\) and \(yf(a')\).

\begin{proposition}\label{v4-arith:gluing:prop:products}
Both \((P,D,\mu_L)\) and \((P,D,\mu_R)\) are differential graded
associative algebras. If \(A,B,C\) and both maps are unital, their common
unit is \(e=(1_A,1_C,0)\). The bilinear map of degree minus one
\[
 K(x,x')=(0,0,(-1)^{p+1}yy')
\]
satisfies
\begin{equation}\label{v4-arith:gluing:eq:homotopy}
 DK+K(D\otimes1+1\otimes D)=\mu_L-\mu_R.
\end{equation}
Consequently the two products induce the same multiplication on
\(H^*(P)\).
\end{proposition}

\begin{proof}
Write \(|x'|=q\). The third component of either bracketing of the left
product on \(x,x',x''\) is
\[
 (-1)^{p+q}f(a)f(a')y''
 +(-1)^pf(a)y'g(c'')+yg(c')g(c'').
\]
For the right product the common third component is
\[
 yf(a')f(a'')+(-1)^pg(c)y'f(a'')
                 +(-1)^{p+q}g(c)g(c')y''.
\]
The other components are associative because \(A\) and \(C\) are.
In the left Leibniz identity the two mixed terms \(f(a)g(c')\)
cancel; the remaining terms give the differential of the third
component in \eqref{v4-arith:gluing:eq:products}. In the right identity
the cancelling terms are \(g(c)f(a')\). This proves compatibility with
\(D\). Substitution proves both unit identities.

The third components of the three terms on the left of
\eqref{v4-arith:gluing:eq:homotopy} are
\[
 \begin{split}
 &(-1)^p(d_By)y'-y(d_By'),\\
 &(-1)^p\bigl(f(a)-g(c)-d_By\bigr)y',\\
 &-y\bigl(f(a')-g(c')-d_By'\bigr).
 \end{split}
\]
Their sum is the difference of the two products. On cocycles this
difference is the boundary \(DK(x,x')\).
\end{proof}

The left product also has a finite interval model. Let \(I_R\) have
basis \(e_0,e_1\) in degree zero and \(h\) in degree one, with
\[
 de_0=-h,\quad de_1=h,\quad dh=0,
 \qquad e_ie_j=\delta_{ij}e_i,
 \qquad e_0h=h=he_1,
\]
and \(e_1h=he_0=h^2=0\). Its unit is \(e_0+e_1\). The evaluations
\(\operatorname{ev}_i:I_R\to R\) send \(e_i\) to one and the other
basis elements to zero. The graded algebra is the path algebra of one
arrow from \(0\) to \(1\), placed in degree one. Its differential is
the graded commutator with \(h\), so Leibniz and \(d^2=0\) follow
from \(h^2=0\). The constant inclusion \(R\to I_R\)
is a quasi-isomorphism: after quotienting by \(R(e_0+e_1)\), the
differential identifies two free rank-one modules.

\begin{proposition}\label{v4-arith:gluing:prop:interval}
The map
\[
 (a,c,y)\longmapsto
 \bigl(a,c,f(a)e_0+g(c)e_1+(-1)^p yh\bigr),\qquad |(a,c,y)|=p,
\]
is a strict isomorphism from \((P,D,\mu_L)\) onto the dg algebra
\[
 \left\{(a,c,\beta)\in A\times C\times(B\otimes_R I_R):
       \operatorname{ev}_0\beta=f(a),\quad
       \operatorname{ev}_1\beta=g(c)\right\}.
\]
Here evaluations act on the interval factor, and \(be_i\) and \(bh\)
denote tensors in \(B\otimes_R I_R\).
\end{proposition}

\begin{proof}
Every triple in the target has the indicated unique form. Its edge
coefficient after differentiation is
\((-1)^{p+1}(f(a)-g(c)-d_By)\).
In a product of degrees \(p,q\), the edge coefficient is
\[
 (-1)^qf(a)y'+(-1)^{p+q}yg(c').
\]
This is \((-1)^{p+q}\) times the third component of \(\mu_L\).
\end{proof}

\section{The integral bar comparison}
\label{v4-arith:gluing:sec:bar}

Let \(s:P\to P[1]\) have degree minus one, and write \(W=sP\).
Thus \(|sx|=|x|-1\); the third coordinate of
\(s(a,c,y)\) identifies with \(y\in B\) in this degree. The tensor
coalgebra is
\[
 T_R^c(W)=\bigoplus_{n\geq0}W^{\otimes_R n}.
\]
Its coproduct \(\Delta\) cuts a word at every position, including the
two ends. Its coaugmentation \(\eta:R\to T_R^c(W)\) includes the empty
word, and its counit \(\epsilon\) projects onto that word. On the
kernel of \(\epsilon\), removing the two end cuts gives the reduced
coproduct. The property that a sufficiently long iteration kills each
element of this kernel is called
\emph{conilpotence}; the tensor coalgebra has this property.

Define \(\mathcal B_L\) and \(\mathcal B_R\) to be this coalgebra with
degree-one coderivations determined by
\begin{equation}\label{v4-arith:gluing:eq:bar-signs}
 b_1(sx)=-sDx,\qquad
 b_{2,L/R}(sx,sx')=(-1)^{|x|}s\mu_{L/R}(x,x').
\end{equation}
Insertion of either component after \(w_1,\ldots,w_j\) has sign
\((-1)^{|w_1|+\cdots+|w_j|}\). All other components vanish, and the
differential of the empty word is zero. The dg algebra identities imply
\(b_L^2=b_R^2=0\).

These are the bars of the underlying nonunital algebras. Equivalently,
they are the reduced bars of their unitizations \(U_{L/R}=R\oplus P\),
with multiplication
\[
 (r,p)(r',p')=(rr',rp'+r'p+\mu_{L/R}(p,p'))
\]
and augmentation given by projection onto \(R\). No augmentation of \(P\) is chosen.
No flatness or boundedness assumption is needed for these coalgebras.

\begin{theorem}[Integral endpoint comparison]
\label{v4-arith:gluing:thm:comparison}
There are inverse dg coalgebra isomorphisms
\[
 F:\mathcal B_L\longrightarrow\mathcal B_R,
 \qquad G:\mathcal B_R\longrightarrow\mathcal B_L.
\]
Their components with output in \(W\) are
\begin{equation}\label{v4-arith:gluing:eq:taylor}
 \begin{split}
 F_1&=G_1=1_W,\\
 F_n(w_1,\ldots,w_n)&=s(0,0,y_1\cdots y_n),\qquad n\geq2,\\
 G_n&=(-1)^{n-1}F_n,\qquad n\geq2,
 \end{split}
\end{equation}
where \(w_i=s(a_i,c_i,y_i)\). Both maps preserve the coaugmentation,
counit and increasing filtration by tensor length. They commute with
strict maps of cospans and ordinary scalar extension. If the cospan is
unital, every higher component vanishes when an input is \(se\).
In particular, the two endpoint algebras are strictly unital
\(A_\infty\)-isomorphic in that case.
\end{theorem}

\begin{proof}
On an input of length \(N\), define \(F\) by summing over all partitions
into consecutive nonempty blocks. Apply \(F_i\) to a block of length
\(i\), and tensor the block values in their original order. Set
\(F(1)=1\). Cutting an output between blocks is equivalent to cutting
the input there and partitioning its two pieces. This proves the
coalgebra identity. Every sum is finite, and no output is longer than
the input.

For the chain identity, set \(|w_i|=q_i\) and
\(Q_j=q_1+\cdots+q_j\), with \(Q_0=0\). Write \(b_1^{\otimes}\)
for the tensor differential on an input word. Since
\[
 \bigl(b_1s(a,c,y)\bigr)_B=d_By-f(a)+g(c),
\]
the differential of the ordered product cancels its internal
\(d_By_i\) terms against \(F_nb_1^{\otimes}\). For \(n\geq2\), the
third component of \(b_1F_n-F_nb_1^{\otimes}\) is
\begin{equation}\label{v4-arith:gluing:eq:unary-defect}
 \sum_{j=1}^n(-1)^{Q_{j-1}}
 y_1\cdots y_{j-1}\bigl(f(a_j)-g(c_j)\bigr)y_{j+1}\cdots y_n.
\end{equation}
At arity two,
\[
 \begin{split}
 (b_{2,L})_B&=f(a_1)y_2-(-1)^{q_1}y_1g(c_2),\\
 (b_{2,R})_B&=-(-1)^{q_1}y_1f(a_2)+g(c_1)y_2.
 \end{split}
\]
Their difference is exactly \eqref{v4-arith:gluing:eq:unary-defect}.
At arity \(n\geq3\), an interior split in
\(b_{2,R}(F_i,F_{n-i})\) has two purely local outputs and is zero.
The two end splits give
\[
 g(c_1)y_2\cdots y_n
       -(-1)^{Q_{n-1}}y_1\cdots y_{n-1}f(a_n).
\]
Adding these to \eqref{v4-arith:gluing:eq:unary-defect} leaves
\[
 \begin{split}
 &\sum_{j=1}^{n-1}(-1)^{Q_{j-1}}
       y_1\cdots y_{j-1}f(a_j)y_{j+1}\cdots y_n\\
 &\quad-\sum_{j=2}^{n}(-1)^{Q_{j-1}}
       y_1\cdots y_{j-1}g(c_j)y_{j+1}\cdots y_n.
 \end{split}
\]
This is the third component of the source binary contribution
\(F_{n-1}b_{2,L}\), including its insertion signs. The other components
vanish in these arities. Thus \(b_RF=Fb_L\) after projection to \(W\).
The two sides are coderivations relative to \(F\), so this projection
determines their equality. Explicitly, induction on input length makes
their remaining difference primitive; a primitive of the tensor
coalgebra with zero projection to \(W\) is zero.

Construct \(G\) by the same partition formula. On length \(n\geq2\),
the component in \(W\) of \(FG\) is the local product
\(s(0,0,y_1\cdots y_n)\) times
\[
 \sum_{k=1}^n\binom{n-1}{k-1}(-1)^{n-k}=0.
\]
For \(GF\) the sign is \((-1)^{k-1}\), and the sum again vanishes.
The linear components are identities. The uniqueness of coalgebra maps
with prescribed components now gives \(FG=GF=1\); hence \(G\) is also a
chain map.

A strict map of cospans takes an ordered product of the \(y_i\) to the
ordered product of their images. Ordinary scalar extension preserves
the same formulas and finite sums. Finally the local coordinate of
\(e\) is zero. This proves the unit assertion.
\end{proof}

The comparison contains more than the binary homotopy
\eqref{v4-arith:gluing:eq:homotopy}. For example, its arity-three
component is \(F_3(w_1,w_2,w_3)=s(0,0,y_1y_2y_3)\).
Let \(\tau_{L/R}:\mathcal B_{L/R}\to P\) be the degree-one map
that desuspends the length-one part and kills every other length.
Then \(\tau_RF\) evaluates every higher product in
\eqref{v4-arith:gluing:eq:taylor}. Those terms carry the change of
endpoint into representations.

\section{Transporting all module operations}
\label{v4-arith:gluing:sec:modules}

For a graded \(R\)-module \(M\), the tensor product
\(M\otimes_R T_R^c(W)\) is a right comodule by deconcatenation on the
word factor. Denote its coaction by \(\rho\). A right \(A_\infty\)
module over \((P,\mu_R)\) is a degree-one map \(Q_M\) satisfying
\[
 \rho Q_M=(Q_M\otimes1+1\otimes b_R)\rho,
 \qquad Q_M^2=0.
\]
It is determined by its components
\(\beta_n:M\otimes_R W^{\otimes_R n}\to M\) of degree one, including
the differential \(\beta_0\). On a word, \(Q_M\) sums the operations
\(\beta_j\) on the module entry and the first \(j\) factors, together
with \((-1)^{|m|}m\otimes b_R\) on the word. All these sums are finite.
The graded maps in the module category are the graded comodule maps
\(T:M\otimes\mathcal B_R\to N\otimes\mathcal B_R\), with
differential \(Q_NT-(-1)^{|T|}TQ_M\).

In these unsuspended-module conventions a strict dg right action is
encoded by
\begin{equation}\label{v4-arith:gluing:eq:module-sign}
 \beta_0(m)=d_Mm,\qquad
 \beta_1(m;sx)=(-1)^{|m|+1}mx,\qquad \beta_n=0\quad(n>1).
\end{equation}
The minus sign is fixed by \eqref{v4-arith:gluing:eq:bar-signs}:
on \(m\otimes sx\otimes sx'\), the two action terms in \(Q_M^2\)
have opposite coefficients \((-1)^{|x|}\) and
\((-1)^{|x|+1}\). On one bar factor the identity is the dg module
Leibniz rule.

Strict unitality means \(\beta_1(m;se)=(-1)^{|m|+1}m\) and
vanishing of \(\beta_n\), for \(n>1\), when a word input is \(se\).
A normalized map has the analogous vanishing in every positive
word arity.

\begin{theorem}\label{v4-arith:gluing:thm:modules}
Pullback along \(F\) gives an isomorphism of the dg categories of
right \(A_\infty\) modules over \((P,\mu_R)\) and \((P,\mu_L)\),
with inverse pullback along \(G\). It leaves the underlying complex
unchanged. Its positive-arity operations are
\begin{equation}\label{v4-arith:gluing:eq:module-pullback}
 \beta_n^F(m;w_1,\ldots,w_n)=
 \sum_{\substack{i_1+\cdots+i_k=n\\i_j>0}}
 \beta_k\bigl(m;F_{i_1}(w_1,\ldots),\ldots,
                         F_{i_k}(\ldots,w_n)\bigr).
\end{equation}
For unital cospans it preserves strict units and normalized module maps.
\end{theorem}

\begin{proof}
Let \(C_F^M=1_M\otimes F\) and \(C_G^M=1_M\otimes G\). Define
\[
 Q_M^F=C_G^M Q_M C_F^M.
\]
The identity \(b_RF=Fb_L\) transports the comodule identity to one
over \(b_L\), and \((Q_M^F)^2=C_G^M Q_M^2 C_F^M=0\).
Projection to the length-zero module factor satisfies
\(\operatorname{pr}_M C_G^M=\operatorname{pr}_M\). Its composite with
\(Q_M^F\) is therefore exactly
\eqref{v4-arith:gluing:eq:module-pullback}. No extra signs occur because
every \(F_i\) has degree zero.

A homogeneous comodule map \(T:M\otimes\mathcal B_R\to
N\otimes\mathcal B_R\) is sent to \(C_G^NTC_F^M\). Conjugation
preserves composition and the differential on maps
\(Q_NT-(-1)^{|T|}TQ_M\). Conjugation by the inverse maps reverses
both constructions. This proves the assertion for the entire dg
category.

A unit input in a surviving partition must be a singleton,
since every higher \(F_i\) vanishes on it. The unit identities for
the outer operation then give exactly these properties after pullback.
The same argument applies to \(G\).
\end{proof}

A strict action can acquire operations in every arity:
\(\beta_n^F(m;w_1,\ldots,w_n)=\beta_1(m;F_n(w_1,\ldots,w_n))\).
The identity linear term of \(F\) therefore does not permit the module
comparison to be replaced by the identity on its underlying complex.

\section{Relative bimodules}
\label{v4-arith:gluing:sec:bimodules}

An interface has actions on both sides. To compare them, fix two cospans
\[
 A_0\xrightarrow{f_A}A_1\xleftarrow{g_A}A_2,
 \qquad E_0\xrightarrow{f_E}E_1\xleftarrow{g_E}E_2.
\]
Let \(M_i\) be a dg \(A_i\)--\(E_i\) bimodule, and choose chain maps
\(\phi:M_0\to M_1\), \(\psi:M_2\to M_1\) with
\[
 \phi(ame)=f_A(a)\phi(m)f_E(e),\qquad
 \psi(ame)=g_A(a)\psi(m)g_E(e)
\]
on the respective domains. Write \(P_A,P_E\) for their relative
algebra complexes, and put
\[
 P_M^n=M_0^n\oplus M_2^n\oplus M_1^{n-1},\qquad
 D_M(m_0,m_2,v)=(dm_0,dm_2,\phi(m_0)-\psi(m_2)-dv).
\]
For \(x=(a_0,a_2,y)\in P_A^p\), \(m=(m_0,m_2,v)\in P_M^q\),
and \(u=(u_0,u_2,z)\in P_E\), the endpoint actions are
\begin{align*}
 x\cdot_Lm&=(a_0m_0,a_2m_2,(-1)^pf_A(a_0)v+y\psi(m_2)),\\
 x\cdot_Rm&=(a_0m_0,a_2m_2,y\phi(m_0)+(-1)^pg_A(a_2)v),\\
 m\cdot_Lu&=(m_0u_0,m_2u_2,(-1)^q\phi(m_0)z+vg_E(u_2)),\\
 m\cdot_Ru&=(m_0u_0,m_2u_2,vf_E(u_0)+(-1)^q\psi(m_2)z).
\end{align*}
For each choice of endpoint these make \(P_M\) a dg bimodule for the
matching products on \(P_A\) and \(P_E\). The proof below includes
the compatibility of the two actions.

The corresponding bar bicomodule is
\[
 \mathcal B^M_{L/R}
 =T_R^c(sP_A)\otimes_R sP_M\otimes_R T_R^c(sP_E).
\]
Its left and right coactions cut before or after the middle entry.
Its differential inserts the unary differentials, the algebra
products, and the two actions, with the suspended conventions of
\eqref{v4-arith:gluing:eq:bar-signs} on every entry, including \(sP_M\).
In particular this middle entry is suspended, unlike the module entry
in \eqref{v4-arith:gluing:eq:module-sign}.

\begin{theorem}\label{v4-arith:gluing:thm:bimodules}
There is a dg bicomodule isomorphism
\(\Phi:\mathcal B^M_L\to\mathcal B^M_R\) over the coalgebra maps
\(F_A,F_E\) of Theorem~\ref{v4-arith:gluing:thm:comparison}. Its
component with no algebra inputs is the identity on \(sP_M\). For
\(r+t\geq1\), its component with \(r\) left and \(t\) right inputs is
\begin{equation}\label{v4-arith:gluing:eq:bimodule-taylor}
 \Phi_{r,t}(sx_1,\ldots,sx_r;sm;su_1,\ldots,su_t)
 =s(0,0,y_1\cdots y_r\,v\,z_1\cdots z_t).
\end{equation}
Its inverse has higher components
\(\Psi_{r,t}=(-1)^{r+t}\Phi_{r,t}\) and identity linear term.
The maps commute with strict diagram maps and ordinary scalar
extension. They preserve strict units in the algebra inputs when
the cospans and bimodules are unital.
\end{theorem}

\begin{proof}
Form the triangular dg algebras
\[
 T_i=\begin{pmatrix}A_i&M_i\\0&E_i\end{pmatrix},\qquad
 \begin{pmatrix}a&m\\0&u\end{pmatrix}
 \begin{pmatrix}a'&m'\\0&u'\end{pmatrix}
 =\begin{pmatrix}aa'&am'+mu'\\0&uu'\end{pmatrix}.
\]
Their entrywise differential has the dg algebra Leibniz property
because the two module actions do. The maps \(\phi,\psi\), together
with the algebra maps, give a strict cospan \(T_0\to T_1\leftarrow T_2\).
Its relative complex has diagonal components \(P_A,P_E\) and
off-diagonal component \(P_M\). The two endpoint products give
exactly the four displayed actions. Their dg bimodule identities
are thus the corresponding entries of the dg algebra identities.

In the tensor bar of this triangular relative algebra, take words
with upper-left entries, then one upper-right entry, then lower-right
entries. They form \(\mathcal B^M_{L/R}\), and its differential
preserves this submodule. Applying the partition formula for the
triangular algebra preserves it too. The unique block containing
the upper-right entry has local component
\(y_1\cdots y_r v z_1\cdots z_t\), without any interchange of
factors. This proves \eqref{v4-arith:gluing:eq:bimodule-taylor}.
The inverse sign is \((-1)^{(r+t+1)-1}\).

The chain and inverse identities restrict from those of the
triangular bars. Deconcatenation gives
\[
 \Delta_L\Phi=(F_A\otimes\Phi)\Delta_L,
 \qquad\Delta_R\Phi=(\Phi\otimes F_E)\Delta_R.
\]
Finally, strict diagram maps preserve every ordered local product,
scalar extension preserves the finite formulas, and a gluing unit
has local component zero. These give the remaining assertions.
\end{proof}

An algebra local condition for arithmetic cochains must provide a dg
algebra and a multiplicative map into the local cochains. A linear local
condition may instead enter the bimodule construction once its two
actions and the compatible maps are supplied. Its cone as a complex
does not determine either action.

\section{The homological meaning of the bar}
\label{v4-arith:gluing:sec:derived}

The relative complex and its tensor bar have different meanings.
The complex \(P\) computes the relative cohomology of the cospan. The
bar computes a derived tensor of the augmentation module of
\(U=R\oplus P\) under the \(K\)-flatness hypothesis of
Proposition~\ref{v4-prop:bar-derived-tensor}.

\begin{proposition}\label{v4-arith:gluing:prop:derived}
If \(P\) is \(K\)-flat over \(R\), then for either endpoint product
its bar represents \(R\otimes_U^{\mathbf L}R\), where \(R\) is a
left or right \(U\)-module through the augmentation.
\end{proposition}

\begin{proof}
For either endpoint product, the augmentation ideal of
\(U=R\oplus P\) is \(P\), and the differential and bar signs are
those of \eqref{v4-eq:bar-coefficient-signs}. Apply
Proposition~\ref{v4-prop:bar-derived-tensor} to this augmented algebra.
\end{proof}

Degreewise flatness does not replace \(K\)-flatness for an unbounded
complex. Without that hypothesis the explicit bar and the endpoint
isomorphism still exist; the proposition does not identify that bar
with a derived tensor.

The unital gluing case gives a useful distinction even without
flatness. If \(P\) has a closed unit \(e\), its underlying nonunital
bar contracts onto the empty word. On a nonempty word \(c\), set
\(h(c)=se\otimes c\), and set \(h(1)=0\). Since
\(b_1(se)=0\) and \(b_2(se,sx)=sx\), direct substitution gives
\[
 bh+hb=1-\eta\epsilon.
\]
Thus \(H^*(\mathcal B)=R\), although \(H^*(P)\) can be much larger.
Algebraically the same distinction appears in
\(U\cong R\times P\), via \((r,p)\mapsto(r,re+p)\).
To form a reduced bar of an augmented unital algebra \(P\) itself,
one instead uses its specified augmentation ideal. No such choice is
implicit in the endpoint construction.

\section{Homotopy-commuting coefficient maps}
\label{v4-arith:gluing:sec:maps}

Strict maps of cospans preserve the preceding constructions. A
comparison commuting only up to cochain homotopy has an additional
product defect, visible already in the relative case \(C=0\).
Write \(P(f)=A\oplus B[-1]\) for this specialization of
\eqref{v4-arith:gluing:eq:complex}. Alongside \(f:A\to B\),
choose a dg algebra map \(f':A'\to B'\).
Let \(\alpha:A\to A'\) and \(\gamma:B\to B'\) be dg algebra maps,
and let \(H:A\to B'\) have degree minus one with
\[
 f'\alpha-\gamma f=dH+Hd.
\]
Define
\(T:P(f)\to P(f')\) by
\(T(a,y)=(\alpha(a),\gamma(y)+H(a))\).

\begin{proposition}\label{v4-arith:gluing:prop:map-defect}
The map \(T\) is a cochain map. For the left endpoint products its
multiplicativity defect has zero first component and local component
\[
 \begin{split}
 &(-1)^p(f'\alpha-\gamma f)(a)\gamma(y')\\
 &\qquad+(-1)^pf'\alpha(a)H(a')-H(aa'),\qquad |(a,y)|=p.
 \end{split}
\]
It induces a strict tensorwise map of the relative bars if and only
if this expression vanishes on every pair of inputs.
\end{proposition}

\begin{proof}
The local component of \(D'T-TD\) is
\((f'\alpha-\gamma f)(a)-dH(a)-H(da)=0\).
Expanding \(\mu_L'(Tx,Tx')-T\mu_L(x,x')\) gives the displayed
formula. Projection of the strict bar-map identity onto one output
factor on two inputs is precisely multiplicativity; on one input it
is the chain-map identity. These two identities suffice in every
length because the bar differentials have only unary and binary
components.
\end{proof}

If the defect is nonzero, higher components must be constructed in
the permitted operation spaces. Closedness of an obstruction in a
larger Hom complex does not provide a primitive with prescribed
support, pole bounds, units or continuity. The endpoint comparison
above supplies such components for its explicit algebraic problem;
it does not supply them for an arbitrary homotopy-commuting square.

\section{Finite precision and continuous gluing}
\label{v4-arith:gluing:sec:completion}

Fix a prime \(\ell\), put \(R_j=\mathbb Z/\ell^j\), and take compatible
strict cospans over the rings \(R_j\). Write \(W_j=sP_j\) and
\(\mathcal B_{L/R,j}=T_{R_j}^c(W_j)\). Define their restricted
inverse limits by
\[
 \widehat{\mathcal B}_{L/R}
       =\varprojlim_j\bigoplus_{n\geq0}W_j^{\otimes_{R_j}n}.
\]
Each finite-precision module is discrete; the limits have the
inverse-limit topology. An element has only finitely many nonzero
word lengths at any fixed precision, although the number of such
lengths may grow with \(j\). The completed tensor recipient is defined
here to be
\begin{equation}\label{v4-arith:gluing:eq:completed-tensor}
 \widehat{\mathcal B}\,\widehat\otimes\,\widehat{\mathcal B}
       :=\varprojlim_j
           (\mathcal B_j\otimes_{R_j}\mathcal B_j).
\end{equation}
This definition makes no identification with a different completed
tensor product.

For a compatible tower of modules or relative bimodules, use
\[
 \widehat C(M)=\varprojlim_j(M_j\otimes_{R_j}\mathcal B_j),
 \qquad \widehat{\mathcal B}^M=\varprojlim_j\mathcal B_j^M,
\]
with the same topology. Their coaction recipients are, respectively,
\(\varprojlim_j(C(M_j)\otimes_{R_j}\mathcal B_j)\) and the two
limits of \(\mathcal B_{A,j}\otimes_{R_j}\mathcal B_j^M\) and
\(\mathcal B_j^M\otimes_{R_j}\mathcal B_{E,j}\). The subscripts
\(A,E\) distinguish the left and right coefficient bars.

\begin{theorem}\label{v4-arith:gluing:thm:completion}
The maps \(F_j,G_j\) induce inverse continuous maps
\(\widehat F:\widehat{\mathcal B}_L\to\widehat{\mathcal B}_R\) and
\(\widehat G\). They preserve the differential and coproduct, with
recipient \eqref{v4-arith:gluing:eq:completed-tensor}. Compatible
module and bimodule bars admit the corresponding continuous
comparisons, with coaction recipients the inverse limits of their
finite-precision tensor recipients.
\end{theorem}

\begin{proof}
At precision \(j\), each element is a finite sum of words, and each
word has finitely many consecutive partitions. The maps are therefore
defined there and satisfy all chain, inverse and coaction identities.
Strict compatibility with coefficient reduction gives the maps on
inverse limits. Every identity holds after each projection, hence in
the limit. Each output at precision \(j\) depends only on the input
at the same precision, proving continuity. The same argument applies
to \(1_M\otimes F_j\), to module conjugation and to the triangular
bicomodule maps.
\end{proof}

No interchange of cohomology with inverse limit occurs in this proof.
Such an interchange would need its own hypotheses on derived inverse
limits. Ordinary conilpotence can also fail in the completed object.
For example, let \(P=\mathbb Z_\ell a\), with \(|a|=1\), zero
differential and zero product. The element
\[
 z=\sum_{n\geq1}\ell^n(sa)^{\otimes n}
\]
exists in the restricted bar. Every iterated reduced coproduct has
a nonzero component, obtained by cutting a sufficiently long word.
Each finite coefficient bar is conilpotent, while \(z\) is not.

\begin{proposition}\label{v4-arith:gluing:prop:noncontinuity}
The map \(F\) need not extend continuously to the unrestricted
product \(\prod_{n\geq0}W^{\otimes_R n}\) with discrete word
coordinates.
\end{proposition}

\begin{proof}
Take the relative identity map \(f:\mathbb Z\to\mathbb Z\), with
\(C=0\), and let \(v=s(0,1)\). It has degree zero and satisfies
\(F_n(v^{\otimes n})=v\) for every \(n\geq1\).
The partial sums \(z_N=\sum_{n=1}^Nv^{\otimes n}\) converge in the
product of word spaces: each coordinate eventually stabilizes.
The length-one coordinate of \(F(z_N)\) is \(Nv\), which does not
converge in the discrete module \(\mathbb Z v\). Continuity of an
extension would contradict this sequence.
\end{proof}

In the restricted \(\ell\)-adic bar the weighted series
\(\sum_{n\geq1}\ell^nv^{\otimes n}\) is defined. Its length-one image
is \(\ell(1-\ell)^{-1}v\). The denominator is an \(\ell\)-adic unit;
it arises by summing the convergent coefficients, not in the definition
of any Taylor component. The comparison is integral at every finite
precision. An analytic realization must separately bound the higher
components in its chosen seminorms and compare its completed tensors
with the recipient in which the coproduct is defined.
